\chapter{Conclusions and Future Work}
\label{chap:conclusions}

\section{Summary}
\label{sec:concl-summary}

AI shopping assistants have moved from research to production at
hyperscale. The published evidence on the largest deployment
documents systematic failures in recommendation accuracy, factual
correctness, and intrusiveness. The question behind this thesis
was whether those failures are inherent to the AI-shopping
problem or whether they come from integration choices a
different system could have made differently. This thesis argues
for the second reading and builds a prototype that shows it.

The prototype combines a two-stage recommendation pipeline
(classical candidate generation followed by LLM rerank and
explanation), a reactive-first advisor surface driven by
lightweight behavior signals, a review-summarization module
grounded in real review text with explicit conflict surfacing,
an AI-assisted comparison flow with no winner verdict, and an
output-validation layer that strips claims the model could
fabricate from free text and checks every cited identifier
against the live catalog.
A custom evaluation harness sits alongside the prototype. It
scores configurations on grounded quality, latency, and cost
across a canonical eight-configuration sweep.

The measured evidence supports the central claim on two of three
axes. Seven of eight configurations stay within a p95 latency of
7.5~s; only \texttt{gpt-oss-120b} at high reasoning effort
rises to 12.1~s. Every configuration costs well under one cent
per query and produces grounded output across the sweep
(\texttt{groundedness} mean $\ge 0.89$ for every configuration).
The harness adds a finding the claim did not anticipate: within
those bounds there is a real design space where the right choice
is the cheaper, lighter configuration, not the one with the
highest raw quality. Higher reasoning
effort and looser step budgets are uniformly worse on composite.
This disproves the implicit ``more reasoning is always better''
assumption that production code tends to encode.

\section{Practical Implications}
\label{sec:concl-implications}

Five recommendations follow from the design and the measured
results.

\begin{description}
  \item[Never let the LLM generate factual product data.]
    Hydration-over-generation reduces the price-hallucination
    surface to zero by construction. Cited identifiers are
    verified against the live catalog; price and specification
    claims are stripped from free text before render.
  \item[Surface review conflicts explicitly with counts.] A
    summary that hides the dissenting minority gives the
    wrong opinion when the minority is large enough to matter;
    the explicit-count format makes that minority visible.
  \item[Keep AI advice user-initiated.] Behavior signals prepare
    context but do not auto-fire the LLM. The reactive surface
    respects user agency and avoids the intrusiveness pattern
    documented for comparable production systems.
  \item[Route by query type.] The harness shows a real
    cost-quality tradeoff between Flash-Lite and
    \texttt{no-fewshot} \texttt{gpt-oss-120b}, and a uniform
    loss from higher reasoning effort. Configuration choices
    should be backed by measured evidence, not intuition.
  \item[Measure every configuration choice.] The few-shot
    contamination finding is the clearest example: the
    production UI hid the failure mode, and only a harness
    running against a broader query distribution surfaced it.
\end{description}

\section{Limitations}
\label{sec:concl-limitations}

The biggest scope limits are the lack of a human-subject
evaluation, the lack of an offline-IR comparison for the
classical recommendation algorithms, the seeded-catalog scale
(about 200 products, 4\,836 AI-generated reviews), and
replication against only a single deployment. The thesis does
not claim to address any of these. It claims that the
engineering work it does report holds up under measurement
within those limits.

\section{Future Work}
\label{sec:concl-future}

A within-subjects user study (perceived usefulness, trust,
intrusiveness, purchase confidence) is the most valuable
extension: the engineering-quality metrics in
Chapter~\ref{chap:system} do not predict UX outcomes, and human
evaluation would close that gap. Adding offline
recommender-quality metrics (Hit~Rate@5, NDCG@10, catalog
coverage, intra-list diversity, popularity-bias Gini) against a
held-out temporal split of the seeded order data is a low-risk
extension that closes the second methodological gap, because it
shares the same data pipeline. Extending the harness to
multi-category catalogs and non-electronics domains would test
whether the theme-fidelity scorer and the comparison flow
generalize beyond the curated dataset. Richer behavior signals
(mouse trajectory, eye tracking, repeated-visit patterns) are
cheap to capture and worth evaluating once human-subject data is
available. Production hardening at 10\,000+ products and traffic
load is documented in \code{docs/scale-considerations.md}.
Turning the eval harness into a CI regression gate with Pareto
comparison against a reference run would let the design
discipline this thesis argues for survive in a real engineering
workflow.

\section{Closing Remarks}
\label{sec:concl-closing}

The thesis argues, end to end, that trustworthy AI shopping
assistance is a problem of integration discipline, not of model
capability. The failure modes documented for the largest
production system can be addressed through structural decisions
a different system can make differently. A working prototype
plus an evaluation harness are enough to show that those
decisions hold up under measurement.
